\chapter*{Start here: from zero}
\addcontentsline{toc}{chapter}{Start here: from zero}
\markboth{Start here}{Start here}

You do not need to have built a circuit, studied physics, or touched a soldering
iron. This book starts with what electricity \emph{is} and ends with a small CPU
you design yourself. What it does use is a little arithmetic with very large and
very small numbers, so this chapter covers that, then shows how the book is laid
out.

\section*{Powers of ten and the prefixes you will see}

Electronics uses numbers from billionths to millions, so we write them with a
prefix. Memorise this table; it appears on every page of the book.

\begin{center}
\begin{tabular}{lll}
\toprule
\thead{Prefix} & \thead{Symbol} & \thead{Multiply by}\\
\midrule
mega  & M     & $10^{6}$ (one million)\\
kilo  & k     & $10^{3}$ (one thousand)\\
milli & m     & $10^{-3}$ (one thousandth)\\
micro & $\mu$ & $10^{-6}$ (one millionth)\\
nano  & n     & $10^{-9}$\\
pico  & p     & $10^{-12}$\\
\bottomrule
\end{tabular}
\end{center}

So $4.7\,\mathrm{k}\Omega = 4700\,\Omega$, and $2.2\,\mu\mathrm{F} = 2.2\times10^{-6}\,\mathrm{F}$.
Mixing prefixes is the commonest source of silly errors: $1\,\mathrm{mA}$ is
$0.001\,\mathrm{A}$, not $1000\,\mathrm{A}$. Convert to base units (volts, amperes, ohms, farads) before you
calculate, and convert back at the end.

\section*{Rearranging a formula}

Almost every calculation here is one formula used three ways. Take Ohm's law
(Chapter~2), $V=IR$. If you want the current, divide both sides by $R$:
$I=V/R$. For $V=5\,\mathrm{V}$ and $R=1\,\mathrm{k}\Omega=1000\,\Omega$ this gives
$I=5/1000=0.005\,\mathrm{A}=5\,\mathrm{mA}$. If you want the resistance, divide by
$I$ instead: $R=V/I$. Whatever you do to one side of an equation, do to the other,
and keep the units attached: volts divided by ohms gives amperes.

\section*{A word about safety}

Everything in the first half of the book runs on batteries and USB power, and every
chapter says clearly what it will and will not ask you to touch. Mains electricity
(the wall socket) can kill, and nothing in this book asks you to open or modify
anything connected to it. Chapter~1 ends with the full safety statement; read it
before you pick up a tool.

\section*{How the books of this series fit together}

\begin{center}
\begin{tabular}{P{2.4cm}P{9.0cm}}
\toprule
\thead{Book} & \thead{Connection}\\
\midrule
8 Physics for engineers & Explains the \emph{why} behind the analogue chapters (fields, capacitance, induction); its ``In engineering practice'' sections point back here.\\
7 Systems & Picks up exactly where Chapter~16 here ends: the custom 8-bit CPU is mapped to RISC-V and extended to a pipelined core.\\
2 Discrete maths & The logic of Chapter~8 here (Boolean algebra) is the same as the logic in Book~2, Chapter~1.\\
\bottomrule
\end{tabular}
\end{center}

\section*{How to use the boxes in each chapter}

Each chapter opens with \emph{After this chapter you can}, listing what you will
be able to do. From Chapter~2, a \emph{Before you start} box lists the chapters it
relies on with a one-sentence recap. Chapters 1--7 are analogue (voltages and
currents that vary smoothly); from Chapter~8 on, everything is digital (ones and
zeros). You can start at Chapter~8 if you only care about the CPU, but Chapters~1--3
make the rest easier.
